\section*{Цель работы}
Изучение динамического режима средств измерений.

\section*{Задачи}
\begin{enumerate}
    \item Ознакомиться с лабораторной установкой и собрать схему исследования динамического звена 2-го порядка.
    \item Исследовать динамический режим при ступенчатом изменении входного сигнала (прямоугольные импульсы):
    \begin{itemize}
        \item Определить динамическую погрешность в $6\dots10$ точках на одном полупериоде.
        \item Исследовать зависимость времени установления $t_y$ от частоты собственных колебаний $f_0$.
        \item Исследовать зависимость времени установления $t_y$ от коэффициента демпфирования $\beta$.
    \end{itemize}
    \item Исследовать динамический режим при синусоидальном входном воздействии:
    \begin{itemize}
        \item Определить погрешности в $8\dots10$ точках на одном периоде сигнала.
        \item Оценить максимальное (амплитудное) значение динамической погрешности.
    \end{itemize}
\end{enumerate}

\section*{Теоретические сведения}
Динамический режим характеризуется изменением входного сигнала во времени, что влияет на результат измерения из-за инерционных свойств средства измерений (СИ).
Идеальные безынерционные СИ описываются линейной зависимостью:
$$y_u(t) = k_H x(t)$$
где $y_u(t)$ — выходной сигнал идеального СИ, $x(t)$ — входной сигнал, $k_H$ — номинальный коэффициент преобразования.

Реальные средства измерений описываются дифференциальными уравнениями, и их выходной сигнал $y(t)$ отличается от идеального. Разность между реальным и идеальным выходными сигналами при одном и том же входном воздействии определяет динамическую погрешность:
\begin{equation}
    \Delta y(t) = y(t) - y_u(t)
    \label{eq:dynamic_error}
\end{equation}
В данной работе исследуется звено 2-го порядка (фильтр нижних частот), параметры которого (частота собственных колебаний $f_0$ и коэффициент демпфирования $\beta$) могут изменяться дискретно.

\begin{figure}[H]
    \centering
    \includegraphics[width=0.8\linewidth]{figs/scheme_block.pdf}
    \caption{Структурная схема лабораторной установки}
    \label{fig:block_scheme}
\end{figure}

Измерения мгновенных значений сигналов производятся методом стробоскопического отбора с использованием устройств выборки и хранения (УВХ). Момент времени измерения $t_i$ задается задержкой $\tau$ относительно импульса синхронизации.

\section*{Инструкция по выполнению работы}

\subsection*{Подготовка к работе}
\begin{enumerate}
    \item Соберите схему соединений согласно \cref{fig:block_scheme}.
    \begin{itemize}
        \item Соедините выходы генератора сигналов (ГС) и фильтра (ФНЧ) со входами Y1 и Y2 осциллографа соответственно.
        \item Подключите выход «Синхронизация» ко входу внешней синхронизации осциллографа.
        \item Подключите выход импульса подсветки (строб) ко входу Z осциллографа (для визуализации точки выборки).
        \item Соедините выходы УВХ1 и УВХ2 с цифровыми вольтметрами ЦВ1 и ЦВ2.
        \item Подключите выход блока управления выборкой к цифровому частотомеру (ЦЧ) для измерения временной задержки.
    \end{itemize}
    \item Установите на генераторе ГС форму сигнала и частоту согласно заданию преподавателя.
    \item Получите у преподавателя параметры исследуемого звена ($f_0$ и $\beta$).
\end{enumerate}

\subsection*{Часть 1. Исследование переходного процесса (прямоугольный импульс)}
\begin{enumerate}
    \item Установите на ГС режим генерации прямоугольных импульсов.
    \item Установите заданные переключателями на модуле ФНЧ значения $f_0$ и $\beta$.
    \item \textit{Снятие переходной характеристики:}
    \begin{itemize}
        \item Вращая ручку «$\tau = \text{var}$», перемещайте строб-импульс по экрану осциллографа вдоль полупериода сигнала.
        \item Выберите $6\dots10$ точек $t_i$, охватывающих характерные участки переходного процесса (фронт, выброс, установление).
        \item Для каждой точки зафиксируйте время $t_i$ (по частотомеру ЦЧ), входное напряжение $u_{\text{вх} i}$ (по ЦВ1) и выходное напряжение $u_{\text{вых} i}$ (по ЦВ2).
        \item Данные занесите в \cref{tab:step_response}.
    \end{itemize}
    \item \textit{Определение времени установления от частоты $f_0$:}
    \begin{itemize}
        \item Зафиксируйте значение $\beta$ (по указанию преподавателя).
        \item Последовательно устанавливая 4 различных значения $f_0$, определите время установления $t_y$.
        \item Критерий установления: выходной сигнал входит в трубку $\pm 5\%$ от установившегося значения и не выходит из неё.
        \item Результаты занесите в \cref{tab:settling_freq}.
    \end{itemize}
    \item \textit{Определение времени установления от демпфирования $\beta$:}
    \begin{itemize}
        \item Зафиксируйте значение $f_0$.
        \item Последовательно устанавливая 4 значения $\beta$, определите время установления $t_y$.
        \item Результаты занесите в \cref{tab:settling_beta}.
    \end{itemize}
\end{enumerate}

\subsection*{Часть 2. Исследование синусоидального воздействия}
\begin{enumerate}
    \item Переключите ГС в режим синусоидального сигнала.
    \item Установите параметры $f_0$ и $\beta$ согласно заданию.
    \item Выберите $8\dots10$ точек $t_i$, равномерно распределенных по одному периоду колебания.
    \item В каждой точке измерьте $t_i$, $u_{\text{вх} i}$ и $u_{\text{вых} i}$.
    \item Данные занесите в \cref{tab:sine_response}.
\end{enumerate}

\section*{Инструкция по обработке результатов}
\begin{enumerate}
    \item \textit{Для переходного процесса:}
    \begin{itemize}
        \item Постройте на одном графике зависимости $u_{\text{вх}}(t)$ и $u_{\text{вых}}(t)$.
        \item Рассчитайте динамическую погрешность в каждой точке: $\Delta y(t_i) = u_{\text{вых} i} - k_H \cdot u_{\text{вх} i}$ (принимая $k_H = 1$, если не указано иное).
        \item Постройте график динамической погрешности $\Delta y(t)$.
        \item Постройте графики зависимостей $t_y = F(f_0)$ и $t_y = F(\beta)$.
    \end{itemize}
    \item \textit{Для синусоидального воздействия:}
    \begin{itemize}
        \item Постройте совмещенные графики входного и выходного сигналов.
        \item Рассчитайте и постройте график динамической погрешности.
        \item Определите максимальное значение динамической погрешности (амплитуду погрешности).
    \end{itemize}
\end{enumerate}